%! Exponential Growth and Decay — Unit 7, Lesson 7.1 — Warm-Up (Answer Key)
\documentclass[10pt]{article}
\usepackage{algebra2-article}
\usepackage{algebra2-key}   % pulls in -boxes; do NOT also load -boxes
\newcommand{\TallMath}[1]{$\displaystyle #1\rule[-1.4em]{0pt}{3.2em}$}

\begin{document}

\pageheader{Unit 7, Lesson 7.1}{Warm-Up --- Answer Key}
\namedateperiod

\vspace{0.10in}

\begin{notesbox}{Getting Ready: Skills We'll Use Today}
\small These three skills power today's lesson on \emph{building} an exponential model. Work quickly.

\vspace{4pt}
\textbf{1. Percent change --- the long way, then the short way.} A jacket costs \$$50$.

\vspace{3pt}
(a) The price goes \textbf{up} $8\%$. \\[2pt]
\hspace*{1.2em} The increase is \ans{\$4}, so the new price is \ans{\$54}. \\[2pt]
\hspace*{1.2em} Now get the same answer with \emph{one} multiplication:
$\$50 \times$ \ans{1.08} $=$ \ans{\$54}. \\[2pt]
\hspace*{1.2em} The new price is \ans{108}\% of the old price.

\vspace{4pt}
(b) Starting over from \$$50$, the price goes \textbf{down} $15\%$. \\[2pt]
\hspace*{1.2em} The decrease is \ans{\$7.50}, so the new price is \ans{\$42.50}. \\[2pt]
\hspace*{1.2em} One multiplication: $\$50 \times$ \ans{0.85} $=$ \ans{\$42.50}. \\[2pt]
\hspace*{1.2em} The new price is \ans{85}\% of the old price.

\vspace{5pt}
\textbf{2. Reading $a$ and $b$ (Lesson 7.0).} Fill in the table from the equation alone.
\begin{center}
\renewcommand{\arraystretch}{1.5}
\begin{tabularx}{\linewidth}{@{}l X X X X@{}}
\hline
\textbf{Function} & \textbf{$a$} & \textbf{$b$} & \textbf{Growth or decay?} & \textbf{$y$-intercept} \\
\hline
$y=5(1.2)^{x}$ & \ans{5} & \ans{1.2} & \ans{growth} & \ans{$(0,5)$} \\
$y=80(0.6)^{x}$ & \ans{80} & \ans{0.6} & \ans{decay} & \ans{$(0,80)$} \\
\hline
\end{tabularx}
\end{center}

\vspace{2pt}
\textbf{3. A constant ratio.} The table below is exponential.
\begin{center}
\renewcommand{\arraystretch}{1.3}
\begin{tabular}{|c|c|c|c|c|}
\hline
$x$ & $0$ & $1$ & $2$ & $3$ \\ \hline
$y$ & $7$ & $21$ & $63$ & $189$ \\ \hline
\end{tabular}
\end{center}
Divide to find the constant ratio: $\dfrac{21}{7}=$ \ans{3}, \; $\dfrac{63}{21}=$
\ans{3}, \; $\dfrac{189}{63}=$ \ans{3}. \\[3pt]
So the base is $b=$ \ans{3}, and the output at $x=0$ is $a=$ \ans{7}.

\end{notesbox}

\begin{teachernote}
\textbf{Item 1 is the entire lesson in miniature} --- do not rush the debrief. The two answers that
matter are $1.08$ and $0.85$, and the sentence to force out loud is: \emph{after an $8\%$ increase you
have $108\%$ of what you started with, and after a $15\%$ decrease you have $85\%$.} That is why the
base is $1+r$ or $1-r$ and never the percent itself. The ``$\%$ of the old price'' line is the bridge
--- students who can say ``$108\%$'' will write $b=1.08$; students who only see ``$8\%$'' will write
$b=0.08$, today's signature error.

\textbf{Item 2} is a straight recall check from Lesson 7.0. Anyone who cannot pull $a$ and $b$ off the
equation will not be able to reverse the process today; pair them deliberately for the Activity.

\textbf{Item 3} is the other direction of the same idea: the constant ratio \emph{is} $b$, and the
output at $x=0$ \emph{is} $a$. Those two sentences are the whole ``build it from a table'' method, so
name them now --- students should notice they already have both numbers before the notes begin.
\end{teachernote}

\end{document}
